\title{Direct Preference Optimization: Your Language Model is Secretly a Reward Model}

\begin{abstract}
Large unsupervised language models (LMs) are effective at acquiring extensive world knowledge and demonstrating certain reasoning abilities; however, their entirely unsupervised training processes make it challenging to precisely modulate their output. Conventional techniques for enhancing model steerability typically gather human-annotated comparisons of generated responses and subsequently adapt the unsupervised LM to reflect these preferences, most commonly utilizing reinforcement learning from human feedback (RLHF). Yet, RLHF is both intricate and prone to instability, requiring an initial fitting of a reward model that represents human judgment, followed by fine-tuning the LM with reinforcement learning to optimize for this inferred reward while minimizing deviation from the original model’s distribution. This work proposes an alternative reward model parameterization within the RLHF framework that facilitates direct computation of the optimal policy in closed form, enabling resolution of the traditional RLHF task using a straightforward classification loss. Our approach, termed \textit{Direct Preference Optimization} (DPO), delivers stability, strong performance, and computational efficiency, removing the need for sampling from the LM during training or extensive hyperparameter adjustments. Empirical results demonstrate that DPO fine-tunes LMs to align with human preferences at least as effectively as established methods. In particular, DPO outperforms PPO-based RLHF in sentiment control tasks and matches or exceeds response quality benchmarks in summarization and single-turn dialogue, all while offering a significantly simpler implementation and training procedure.
\end{abstract}

\section{Introduction}
Large unsupervised language models (LMs), when trained on massive datasets, demonstrate a range of unexpected abilities~\citep{chowdhery2022palm, brown2020language, touvron2023llama,bubeck2023sparks}. These models are built on data sourced from humans whose aims, preferences, and levels of expertise differ greatly. As a result, not all modeled behaviors are desirable for emulation; for instance, while it may be beneficial for an AI coding assistant to \textit{recognize} frequent programming errors in order to remedy them, we would actually prefer the model, during code synthesis, to emulate the (sometimes infrequent) high-caliber programming skills represented in its corpus. Likewise, it is valuable for a language model to be \textit{informed} of prevalent misconceptions—such as one held by half the population—yet we clearly do not intend for the model to assert such a misconception as fact in half of relevant responses. This illustrates the necessity of separating and \emph{selecting for desired outputs and conduct} from the broader set of a model’s \textit{learned knowledge and skills} to develop AI systems that are both safe and reliably controllable \citep{ouyang2022training}. Standard approaches steer LMs to align with human preferences predominantly through reinforcement learning (RL). Here, we argue that the RL-based objective at the heart of existing approaches can, in fact, be directly optimized by a simple binary cross-entropy objective, significantly streamlining the overall preference learning process.

Generally, these techniques shape language model behavior by leveraging carefully constructed datasets containing human preferences that capture what users consider safe and beneficial. This phase of preference learning follows the initial stage of large-scale unsupervised pre-training on vast text corpora. Although the simplest form of preference learning entails supervised fine-tuning with exemplary human-generated responses, reinforcement learning from human or AI feedback (RLHF/RLAIF; \citep{christiano2017deep,bai2022constitutional}) has emerged as the most effective paradigm. RLHF first trains a reward model based on human preference data, and subsequently employs RL to update the policy such that it generates responses favored by the reward model, while minimizing deviation from the base model. Despite RLHF yielding models of remarkable dialogue and programming proficiency, it comes with substantial complexity relative to supervised learning, requiring multiple LMs and sampling from the LM’s output during training, thereby demanding considerably more computation.

This work introduces a method for optimizing language models to match human preferences directly, without the need for explicit reward modeling or reinforcement learning approaches. We present \textit{{Direct Preference Optimization} (DPO)}, an algorithm that, while implicitly maximizing the same objective as standard RLHF techniques—namely, reward maximization constrained by KL-divergence—offers greater simplicity in both its formulation and practical implementation. Conceptually, DPO updates the model by boosting the log odds of preferred over non-preferred outputs; it further employs a dynamic importance weighting for each training instance, which helps to mitigate the degeneracies observed when naively maximizing probability ratios. Similar to prior approaches, DPO is grounded in a theoretical framework of preference modeling (e.g., the Bradley-Terry model; \cite{bradley1952rankanalysis}), which quantifies the congruence between a reward function and observed preference judgments. However, whereas traditional RLHF pipelines leverage the preference model to first derive a reward signal and subsequently train the policy to maximize this learned reward, DPO reformulates the problem through a variable substitution, allowing the preference loss to be expressed directly in terms of the policy. As a result, DPO enables direct policy optimization from datasets of human preference comparisons, utilizing a straightforward binary cross-entropy loss, and thereby learns the policy corresponding to the optimal implicit reward consistent with human preferences.

The central contribution of this study is Direct Preference Optimization (DPO), an algorithm for preference-based language model training that dispenses with reinforcement learning. Empirical evaluation demonstrates that DPO matches or exceeds the effectiveness of existing strategies, such as RLHF based on PPO, in tasks including sentiment transformation, summarization, and dialogue, for language models with up to 6B parameters.

Self-supervised language models with ever-increasing parameter counts are capable of performing certain tasks in a zero-shot manner \citep{radford2019language} or by leveraging few-shot prompting \citep{gpt3,megatron,chowdhery2022palm}. Nevertheless, adapting these models to downstream tasks and aligning their outputs with user intent is markedly enhanced when they are further trained on corpora comprised of instructions and corresponding human-generated completions \citep{mishra-etal-2022-cross,sanh2022multitask,chung2022scaling,thoppilan2022lamda}. This process, commonly referred to as ‘instruction-tuning,’ facilitates large language models’ ability to generalize beyond the scope of the training instructions and broadly boosts their practical effectiveness \citep{chung2022scaling}. Although instruction-tuning is successful, it is often more feasible to obtain comparative (i.e., relative) human assessments of output quality than to curate gold-standard expert demonstrations. As a result, subsequent research has focused on adapting LLMs using datasets constructed from human preference judgments, leading to notable advancements in translation \citep{kreutzer-etal-2018-reliability}, summarization \citep{stiennon2022learning,ziegler2020finetuning}, story-telling \citep{ziegler2020finetuning}, and following instructions \citep{ouyang2022training,ramamurthy2023is}. These approaches involve first training a neural reward model to reflect the collected preference data—often operationalized via preference modeling techniques such as the Bradley-Terry model \citep{bradley1952rankanalysis}—followed by tuning the language model itself to optimize the learned reward function using reinforcement learning strategies such as REINFORCE \citep{williams1992reinforce}, proximal policy optimization (PPO; \cite{schulman2017proximal}), or related methods \citep{ramamurthy2023is}. Another branch of research employs instruction-tuned LLMs augmented with human feedback to produce synthetic preference data, focusing on specific attributes like safety or harmlessness, with supervision from humans provided mainly through qualitative text rubrics for LLM-guided annotations \citep{bai2022constitutional}. Such methodologies synthesize insights from two main research streams: one that applies reinforcement learning to language model training across various objectives~\citep{Ranzato2015SequenceLT,paulus2018a,wu2018learning}, and another that devises general algorithms for learning from human preferences \citep{christiano2017deep,kupcsik2018learning}. Despite the intuitive benefits of relative preference signals from humans, reinforcement learning-based fine-tuning of large language models is technically demanding; the present work introduces a theoretically-grounded framework for optimizing with relative preferences that avoids the reliance on reinforcement learning.

Beyond language-focused applications, learning policies from preferences has been examined within both bandit and reinforcement learning frameworks, giving rise to a variety of techniques. One notable line of work, contextual dueling bandits (CDB; \cite{yue2012karmed,dudik2015contextual}), addresses situations where actions are chosen based on relative preferences or rankings instead of direct reward signals. In CDB, where absolute rewards are not available, theoretical frameworks introduce the concept of the \textit{von Neumann winner}: a policy guaranteed to achieve at least a 50\% expected success rate against any alternative policy \citep{dudik2015contextual}. Nevertheless, while CDB typically involves the online acquisition of preference information, human preference learning is most often carried out using a static, offline dataset comprising annotated action-preference pairs \citep{yan2022human}. 

In a similar vein, \textit{preference-based RL} (PbRL) focuses on learning from binary preference feedback derived from an underlying, unobserved scoring mechanism, as opposed to explicit reward feedback \citep{BusaFekete2014,ruiz2023dueling}. PbRL encompasses numerous algorithmic strategies, many of which leverage off-policy preference data; these methods generally entail first estimating the hidden reward (scoring) function before subsequently optimizing the policy based on this estimate \citep{jain2013learning,BusaFekete2014,christiano2017deep,sadigh2017active,kupcsik2018learning}. By contrast, we advocate for a direct, single-stage approach that optimizes policies according to observed preferences without an intermediate explicit reward modeling step.

\section{Preliminaries}\label{section:prelims}

Here we summarize the RLHF methodology outlined by \citeauthor{ziegler2020finetuning} (see also \citep{stiennon2022learning, bai2022training, ouyang2022training}). The process is commonly segmented into three main stages: (1) supervised fine-tuning (SFT), (2) collection of preference data and reward modeling, and (3) optimization via reinforcement learning.

\textbf{SFT}: The RLHF procedure typically starts with the supervised fine-tuning of a pre-trained language model, adjusting its parameters with high-quality, task-relevant supervised data (such as for dialogue or summarization tasks), yielding a model $\pi^\text{SFT}$.

\textbf{Reward Modelling Phase}: During the subsequent phase, prompts $x$ are input into the SFT model, which then produces answer pairs $(y_1, y_2)\sim \pi^\text{SFT}(y \mid x)$. Human annotators are then presented with these answer pairs and asked to state a preference—indicating, for example, that $y_w\succ y_l \mid x$, where $y_w$ is the favored answer and $y_l$ the less preferred one from each pair. The assumption is that these preferences result from an underlying, unobserved reward function $r^*(y, x)$. Various techniques exist for modeling these preferences, among which the Bradley-Terry (BT) model \cite{bradley1952rankanalysis} is widely used, although alternatives such as the Plackett-Luce ranking model \citep{plackett1975analysis, luce2012individual} can be incorporated if multiple answers are ranked. Under the BT model, the distribution of human preferences $p^*$ can be formalized as follows:

Given a fixed dataset of pairwise comparisons $\mathcal{D} = \left\{ x^{(i)}, y_w^{(i)}, y_l^{(i)} \right\}_{i=1}^N$ that are drawn according to $p^*$, we introduce a parameterized reward function $r_{\phi}(x, y)$ and learn its parameters by maximizing likelihood. By framing this as a binary classification task, the corresponding negative log-likelihood loss is given by:

\begin{equation}\label{eq:reward_model}
    \mathcal{L}_R(r_{\phi}, \mathcal{D}) = -\mathbb{E}_{(x, y_w, y_l)\sim \mathcal{D}}\left[ \log \sigma \left( r_{\phi}(x, y_w) - r_{\phi}(x, y_l) \right) \right]
\end{equation}

where $\sigma$ denotes the logistic sigmoid. Within large language models (LMs), $r_{\phi}(x, y)$ typically builds upon the SFT model $\pi^\text{SFT}(y \mid x)$, augmented by a linear projection atop the final transformer layer, yielding a single scalar reward prediction~\cite{ziegler2020finetuning}. In order to promote stability in the reward estimates, previous studies have normalized the output, enforcing $\mathbb{E}_{x,y\sim \mathcal{D}}[r_\phi(x, y)] = 0$ for each $x$.

\textbf{RL Fine-Tuning Phase}: Once the reward model is learned, it guides the language model throughout the RL fine-tuning stage. Building on earlier work~\citep{jaques2017sequence, jaques2020human}, policy optimization is defined by the following objective:

\begin{equation}\label{eq:RL}
\max_{\pi_{\theta}} \mathbb{E}_{x\sim \mathcal{D}, y\sim \pi_{\theta}(y \mid x)} \bigl[ r_{\phi}(x, y) \bigr] - \beta \mathbb{D}_{\textrm{KL}}\left[ \pi_{\theta}(y\mid x) \mid \mid \pi_\text{ref}(y\mid x) \right],
\end{equation}

where $\beta$ tunes the strength of the regularization to the reference policy $\pi_\text{ref}$, typically set as the original SFT model $\pi^\text{SFT}$. Practically, $\pi_\theta$ is also initialized from $\pi^\text{SFT}$. This regularization is crucial: it keeps the policy’s outputs within the distribution in which the reward model is reliable, preserves output diversity, and guards against mode collapse towards a handful of high-reward completions. As language is inherently discrete, the optimization objective above is non-differentiable, and thus, reinforcement learning is employed. The prevailing methodology~\citep{ziegler2020finetuning, stiennon2022learning, bai2022training, ouyang2022training} builds a reward function of the form ${r(x, y) = r_{\phi}(x, y) -\beta (\log \pi_{\theta}(y\mid x) - \log \pi_\text{ref}(y\mid x))}$, which is maximized using PPO~\cite{schulman2017proximal}.

\section{Direct Preference Optimization}\label{sec:DPO}

Acknowledging the considerable challenges of scaling RL algorithms for large-scale LM fine-tuning, we aim to develop a more straightforward method for policy optimization directly from preference data. In contrast with earlier RLHF pipelines, which first train a reward model before applying RL, our method exploits a specific reward function form that admits the exact solution for the corresponding optimal policy—thus obviating the need for a reinforcement learning phase. In the following sections, we will elaborate on how we utilize an analytical mapping between reward models and their induced optimal policies. This perspective allows us to translate an objective on rewards to a direct loss on policies by means of a variable substitution. This approach not only circumvents the need to fit an explicit reward model, but also continues to optimize under standard formulations of preference learning (e.g., the Bradley-Terry model). As a result, the policy itself embodies both the language model and the (implicit) reward.

\textbf{DPO Objective Derivation.} Our starting point is the RL objective given in Eq.~\ref{eq:RL}, considering a generic reward function $r$, consistent with prior studies. Building on earlier work~\citep{peters2007reinforcement, peng2019advantage, korbak2022reinforcement, go2023aligning}, the KL-regularized reward optimization in Eq.~\ref{eq:RL} yields an optimal solution expressible as:

\begin{equation}\label{eq:op_policy}
    \pi_r(y\mid x) = \frac{1}{Z(x)}\pi_\text{ref}(y\mid x)\exp\left(\frac{1}{\beta}r(x, y)\right),
\end{equation}

with $Z(x) = \sum_{y} \pi_\text{ref}(y\mid x) \exp\left(\frac{1}{\beta} r(x, y)\right)$ representing the normalizing constant. The complete derivation appears in Appendix~\ref{app:derivation1}. Computing $Z(x)$ remains computationally intractable even when the reward function $r^*$ is approximated by the MLE estimate $r_{\phi}$, as previously noted~\citep{korbak2022reinforcement, go2023aligning}, hindering direct applicability. To address this, we re-express Eq.~\ref{eq:op_policy} by isolating the reward in terms of the policy $\pi_r$, the baseline policy $\pi_\text{ref}$, and the normalizer $Z(\cdot)$. Taking the logarithm of Eq.~\ref{eq:op_policy} and rearranging terms, we arrive at:

\begin{equation}\label{eq:main_eq}
    r(x,y) = \beta \log \frac{\pi_r(y\mid x)}{\pi_\text{ref}(y\mid x)} + \beta \log Z(x).
\end{equation}

This formulation applies equally to the ground-truth reward $r^*$ and its associated optimal policy $\pi^*$. The Bradley-Terry model, importantly, relies on the difference between rewards for a pair of outputs: ${p^*(y_1 \succ y_2 \mid x) = \sigma(r^*(x, y_1) - r^*(x, y_2))}$. Substituting the reparameterization from Eq.~\ref{eq:main_eq} into the preference relation (Eq.~\ref{eq:bradley-terry}) causes the partition terms to cancel, yielding a preference probability that is dependent solely on the optimal policy $\pi^*$ and the reference $\pi_\text{ref}$. The result is that, according to the Bradley-Terry framework, the optimal RLHF policy $\pi^*$ satisfies:

\begin{equation}\label{eq:objective}
    p^*(y_1\succ y_2 \mid x)=\frac{1}{1 + \exp\left(\beta \log \frac{\pi^*(y_2\mid x)}{\pi_\text{ref}(y_2\mid x)} - \beta \log \frac{\pi^*(y_1\mid x)}{\pi_\text{ref}(y_1\mid x)}\right)}
\end{equation}

Details of the derivation appear in Appendix~\ref{app:derivation2}. Although Eq.~\ref{eq:objective} utilizes the Bradley-Terry structure, parallel results can be developed using the broader class of Plackett-Luce models~\citep{plackett1975analysis, luce2012individual}, as outlined in Appendix~\ref{app:plackett_luce_models}.

Having now represented the likelihood of human preference data as a function of the optimal policy rather than the reward function, we can recast the optimization as maximizing the likelihood for a parameterized policy $\pi_\theta$. Mirroring the strategy for reward model objectives (see Eq.~\ref{eq:reward_model}), we write the policy-based loss as:

\begin{equation}\label{eq:optimum_model}
    \mathcal{L}_\text{DPO}(\pi_{\theta}; \pi_\text{ref}) = -\mathbb{E}_{(x, y_w, y_l)\sim \mathcal{D}}\left[\log \sigma \left(\beta \log \frac{\pi_{\theta}(y_w\mid x)}{\pi_\text{ref}(y_w\mid x)} - \beta \log \frac{\pi_{\theta}(y_l\mid x)}{\pi_\text{ref

In this formulation, we infer an implicit reward through a different parameterization, such that the optimal policy directly corresponds to $\pi_\theta$. Notably, because our approach mirrors the structure of a reparameterized Bradley-Terry model, it inherits desirable theoretical characteristics, including consistency, provided certain conditions on the preference data distribution are satisfied \cite{bong2022generalized}. Additional discussion on the theoretical properties of DPO, especially in comparison to related frameworks, is provided in Section~\ref{sec:theory}.

\textbf{How does the DPO update operate?} To gain a mechanistic perspective on DPO, examining the gradient of the loss $\mathcal{L}_\text{DPO}$ is informative. The derivative with respect to the parameter vector $\theta$ can be expressed as follows:

\begin{multline*}\label{eq:gradient}
    \nabla_\theta \mathcal{L}_\text{DPO}(\pi_\theta;\pi_\text{ref}) = \\ -\beta\mathbb{E}_{(x, y_w, y_l) \sim \mathcal{D}} \left[\underbrace{\sigma(\hat{r}_\theta(x, y_l) - \hat{r}_\theta(x, y_w))}_\text{assigns greater emphasis when the reward ordering is incorrect} \left[\underbrace{\nabla_\theta\log \pi(y_w \mid x)}_\text{encourages increasing $y_w$'s probability} - \underbrace{\nabla_\theta\log\pi(y_l \mid x)}_\text{discourages $y_l$'s probability} \right]\right],
\end{multline*}

Here, the function $\hat{r}_\theta(x, y) = \beta \log \frac{\pi_\theta(y \mid x)}{\pi_\text{ref}(y \mid x)}$ represents the implicit reward as derived from the language model $\pi_\theta$ in relation to the baseline model $\pi_\text{ref}$ (with additional discussion in Section~\ref{sec:theory}). In essence, the gradient direction prescribed by $\mathcal{L}_\text{DPO}$ simultaneously amplifies the likelihood of preferred responses $y_w$ and suppresses that of less favored ones $y_l$. The degree to which each training instance influences the update depends on how much the implicit reward model $\hat{r}_\theta$ overestimates the unpreferred option, modulated by the coefficient $\beta$. This weighting quantifies the severity of misordering in the model’s reward assignments and is further shaped by the KL regularization strength. Empirical evidence from our experiments demonstrates the critical role of this weighting scheme, since omitting it—as done in a simpler variant of the method—can lead to undesirable behaviors in the language model, such as degeneration (see Appendix Table~\ref{tab:unlikelihood_generations}).

\textbf{DPO Overview.} The typical DPO workflow consists of the following steps: 1) For each prompt $x$, generate candidate completions $y_1, y_2 \sim \pi_\text{ref}(\cdot \mid x)$, then annotate them with human preferences to form the offline preference dataset $\mathcal{D} = \{x^{(i)}, y_w^{(i)}, y_l^{(i)}\}_{i=1}^N$; and 2) train the language model $\pi_\theta$ by minimizing $\mathcal{L}_\text{DPO}$ using the specified $\pi_\text{ref}$, dataset $\mathcal{D}$, and the desired $\beta$ value. In real-world applications, it is often preferable to leverage publicly available preference datasets rather than creating new samples and collecting additional human annotations. As these datasets are typically derived from $\pi^\text{SFT}$, we set $\pi_\text{ref} = \pi^\text{SFT}$ when such a model is accessible. If $\pi^\text{SFT}$ is not provided, we instead initialize $\pi_\text{ref}$ by maximizing the likelihood over preferred responses ${(x, y_w)}$, that is, ${\pi_\text{ref} = \argmax_{\pi}\mathbb{E}_{x, y_w \sim \mathcal{D}}\left[\log \pi(y_w \mid x)\right]}$. This initialization reduces the mismatch between the unknown true reference distribution and the actual $\pi_\text{ref}$ applied in DPO. Additional implementation specifics and hyperparameter settings are available in Appendix~\ref{app:implementation}.

\section{Theoretical Analysis of DPO} This section aims to further clarify the foundations of the DPO approach, present supporting theoretical analysis, and connect the benefits of DPO to known limitations of actor-critic methods used in RLHF, such as PPO~\cite{schulman2017proximal}.

\label{sec:theory}

\subsection{Your Language Model Is Secretly a Reward Model} DPO streamlines preference optimization by forgoing both explicit reward modeling and reinforcement learning for policy training, instead relying on a unified maximum likelihood objective. In fact, the DPO objective in Eq. \ref{eq:main_eq} corresponds to a Bradley-Terry formulation where the reward is parameterized by $r^*(x, y) = \beta \log\frac{\pi^*_\theta(y \mid x)}{\pi_\text{ref}(y \mid x)}$. The model $\pi_{\theta}$ is then optimized analogously to reward model learning as in Eq. \ref{eq:reward_model}, after a change of variables. The following discussion develops the theoretical underpinnings of this reparameterization, demonstrating both that it imposes no restrictive assumptions on the space of potential reward functions and that it permits exact retrieval of the optimal policy. To establish these claims, we start by introducing an equivalence relation over reward functions.

\begin{definition}
Two reward functions $r(x, y)$ and $r'(x, y)$ are considered equivalent if and only if there exists a function $f$ such that ${r(x, y)-r'(x, y) = f(x)}$.     
\end{definition}

This relation is readily verified to be an equivalence relation, resulting in a partition of the set of all reward functions into equivalence classes. Based on this, we present the following two lemmas:

\begin{lemma}\label{lemma:same_prefrence} Within the Plackett-Luce and, more specifically, the Bradley-Terry framework for modeling preferences, any two reward functions from the same equivalence class generate identical preference distributions.
\end{lemma}

\begin{lemma}\label{lemma:same_policy}
    Any pair of reward functions that belong to the same equivalence class will induce the same optimal policy in the setting of the constrained RL problem.
\end{lemma}

As both results follow directly from the definitions, we provide detailed proofs in Appendix \ref{app:lemma1}. The first lemma highlights a classic issue of model under-specification within the Plackett-Luce family \cite{plackett1975analysis}, necessitating the use of supplementary identifiability restrictions to establish meaningful properties of MLE solutions to Eq. \ref{eq:reward_model} \cite{bong2022generalized}. The second lemma establishes that all reward functions in an equivalence class yield the same set of optimal policies; thus, in our main objective, we need only recover any representative from the optimal equivalence class. In Appendix~\ref{app:thm1}, we establish the following result:

\begin{theorem}\label{thm:main}
    Assuming mild regularity conditions, each reward equivalence class compatible with the Plackett-Luce model (including the Bradley-Terry case) admits a representation of the form ${r(x, y) = \beta \log \frac{\pi(y\mid x)}{\pi_\text{ref}(y\mid x)}}$ for some model $\pi(y\mid x)$ and a given reference policy $\pi_\text{ref}(y \mid x)$.
\end{theorem}

\begin{sproof}
    Take any given reward function $r(x, y)$, which defines an associated optimal policy $\pi_r(y \mid x)$ as outlined in Eq. \ref{eq:op_policy}. We demonstrate that there always exists an equivalent reward function in the same equivalence class as $r$ which fits the specified reparameterized structure. Define the projection operator $f$ as
\begin{equation}
    f(r; \pi_\text{ref}, \beta)(x, y) = r(x, y) - \beta\log\sum_{y}\pi_\text{ref}(y\mid x)\exp\left(\frac{1}{\beta}r(x, y)\right)
\end{equation}
Here, $f$ acts by subtracting the log-partition function with respect to the prefix $x$, effectively normalizing the original reward. Because this normalization term only depends on $x$, $f(r; \pi_\text{ref}, \beta)(x, y)$ remains within the equivalence class of the original reward function $r(x, y)$. Substituting $r$ by the expression in Eq.~\ref{eq:main_eq} (which applies universally for reward functions), it follows that $f(r; \pi_\text{ref}, \beta)(x, y) = \beta \log \frac{\pi_r(y\mid x)}{\pi_\text{ref}(y\mid x)}$. Thus, the mapping $f$ yields a representative of the same equivalence class in the desired canonical form, validating that the suggested reparameterization retains complete generality for our reward model.
\end{sproof}

An alternative interpretation of Theorem~\ref{thm:main} is that it precisely characterizes which reward function, among all equivalent options, is selected by the DPO reparameterization. Specifically, it is the reward function that fulfills:

\begin{equation}\label{eq:lag_p}
     \sum_{y}\underbrace{\pi_\text{ref}(y\mid x)\exp\left(\frac{1}{\beta}r(x, y)\right)}_{=\pi(y\mid x)\text{, using Thm.~\ref{thm:main} reparam.}} = 1,
\end{equation}

In other words, $\pi(y\mid x)$ forms a legitimate probability distribution with all probabilities being non-negative and summing to unity. Furthermore, as per Eq.~\ref{eq:op_policy}, Eq.~\ref{eq:lag_p} can be recognized as the partition function for the optimal policy generated by the reward function $r(x, y)$. The core contribution of the DPO method is the enforcement of constraints within the flexible Plackett-Luce family (especially the Bradley-Terry model) of preference structures. This approach preserves the diversity of possible reward models, while at the same time allowing the optimal policy in Eq.~\ref{eq:op_policy} to remain computationally accessible for every prompt $x$.

\subsection{Instability of Actor-Critic Algorithms} Our analytical framework also offers tools to understand the instability of standard actor-critic approaches, such as PPO, which are common in RLHF pipelines. Focusing on the reinforcement learning fine-tuning phase (see Section~\ref{section:prelims}), we relate this to the control as inference perspective \cite{levine2018reinforcement} applied to the constrained RL task given in \ref{eq:RL}. Here, we consider a parameterized policy $\pi_{\theta}(y\mid x)$ and aim to reduce the divergence $\mathbb{D}_{\text{KL}}[\pi_{\theta}(y|x) \mid \mid \pi^*(y\mid x)]$, where $\pi^*$ refers to the optimal policy from Eq.~\ref{eq:optimum_model} determined by the reward $r_{\phi}(y, x)$. Through further derivation, we obtain the optimization problem:

\begin{equation}\label{eq:AC}
    \max_{\pi_{\theta}}\mathbb{E}_{\pi_{\theta}(y\mid x)}\bigg[\underbrace{r_{\phi}(x, y) -\beta\log\sum_{y}\pi_\text{ref}(y\mid x)\exp\left(\frac{1}{\beta}r_{\phi}(x, y)\right)}_{f(r_{\phi}, \pi_\text{ref}, \beta)} - \underbrace{\beta\log\frac{\pi_{\theta}(y\mid x)}{\pi_\text{ref}(y\mid x)}}_{\text{KL}}\bigg]
\end{equation}

This objective mirrors that which has been optimized in previous studies 
\citep{ziegler2020finetuning, stiennon2022learning, bai2022training, ouyang2022training}, where the reward for the reward class $r_{\phi}$ is equivalent to the DPO framework. In this context, the normalization component within $f(r_{\phi}, \pi_\text{ref}, \beta)$ can be viewed as representing the soft value function for the reference policy $\pi_\text{ref}$. Although this normalization term does not alter the optimal policy, omitting it can cause the policy gradient estimator's variance to increase, potentially destabilizing the learning process. Incorporating the normalization term through a learned value function is feasible, yet optimizing such a function is also challenging. Other approaches have handled normalization by employing a human completion baseline, which corresponds to a single-sample Monte Carlo estimate of the normalization factor. By comparison, the DPO reparameterization produces a reward structure that circumvents the need for any such baselines.

\section{Experiments}
Here, we present an empirical assessment of DPO's effectiveness in policy training directly from preference data. Our initial investigation uses a controlled text-generation task to compare how effectively DPO balances reward maximization against the minimization of KL-divergence from the reference policy, particularly relative to established preference-learning techniques such as PPO. We then extend our analysis to assess DPO's performance when scaling to larger models and more complex RLHF challenges, including tasks in summarization and dialogue. Notably, we observe that DPO generally matches or surpasses robust baselines like PPO-based RLHF and also outperforms strategies that select the best out of $N$ generated trajectories under a trained reward model, all while requiring minimal hyperparameter tuning. We first outline the details of our experimental framework, with further specifics provided in Appendix~\ref{app:exp_details}.

\textbf{Tasks.} We investigate three distinct open-ended text generation tasks in our experimental setup. Across all tasks, the learning algorithms derive a policy based on a dataset of preferences $\mathcal{D}=\bigl\{x^{(i)}, y_w^{(i)}, y_l^{(i)}\bigr\}_{i=1}^N$. For \textbf{controlled sentiment generation}, the input $x$ consists of the initial segment of a movie review sourced from the IMDb dataset \cite{maas-EtAl:2011:ACL-HLT2011}; here, the policy’s objective is to produce an output $y$ with a positive sentiment. To systematically assess performance, we construct preference pairs by generating completions and utilizing a pre-trained sentiment classifier to ensure $p(\text{positive}\mid x,y_w)>p(\text{positive}\mid x,y_l)$. For supervised fine-tuning (SFT), GPT-2-large is further trained to convergence on the training portion of IMDB reviews (additional information can be found in App~\ref{app:sentiment_details}). 

In the \textbf{summarization} task, the input $x$ is a Reddit forum post, with the policy required to summarize its key points as output $y$. Our experiments employ the Reddit TL;DR summarization dataset \citep{volske-etal-2017-tl} alongside human preference data curated by \citeauthor{stiennon2022learning}. For SFT, we use a model adapted to human-authored forum summaries,\footnote{\url{https://huggingface.co/CarperAI/openai_summarize_tldr_sft}} utilizing the TRLX \citep{leandro_von_werra_2023_7790115} framework to perform RLHF. The preference annotations were collected on model outputs sampled from another SFT model that was trained with similar procedures.

The final task, \textbf{single-turn dialogue}, uses $x$ as a human-issued query, ranging from scientific inquiries to personal advice requests. The goal is to generate an informative and engaging reply $y$. This is evaluated using the Anthropic Helpful and Harmless dialogue dataset \citep{bai2022training}, consisting of 170k conversational transcripts between humans and an AI assistant. Each dialogue concludes with a pair of responses from a large (unspecified) language model, with human raters indicating which response they preferred. Since there is no available pre-trained SFT model for this context, we create one by fine-tuning a generic language model exclusively on the preferred responses.

\textbf{Evaluation.} We employ two distinct strategies for evaluating our methods. To assess how well each algorithm optimizes the constrained reward maximization objective in a controlled sentiment generation context, we examine the frontier defined by the reward achieved versus KL-divergence from the reference policy. Since the ground-truth reward function (i.e., the sentiment classifier) is accessible in this controlled setting, it is feasible to compute this frontier precisely. Conversely, in real-world applications, the true reward function is unavailable. Thus, our evaluation relies on measuring the \textit{win rate} of algorithms against a baseline policy. For summary and dialogue tasks, we adopt GPT-4 as a surrogate for human judgment to compare summary quality and response helpfulness, respectively. In the summarization task, the baseline consists of test set reference summaries, whereas in the dialogue task, we use the dataset’s preferred response as the baseline. Prior literature indicates that language models may serve as more reliable automated evaluators compared to conventional metrics \citep{Chen2023ExploringTU}. To validate our reliance on GPT-4 for assessment, we also perform a human evaluation (Sec.~\ref{sec:human-judgments}). Our findings reveal that GPT-4’s judgments align closely with those of human annotators, and human-GPT-4 agreement matches or surpasses inter-annotator agreement levels.

\textbf{Methods.} Besides DPO, our comparative study covers a variety of established techniques for aligning language model outputs with human preferences. For summarization, we utilize zero-shot prompting of \textbf{GPT-J} \citep{gpt-j}, and for the dialogue setting, we use 2-shot prompting with \textbf{Pythia-2.8B} \citep{biderman2023pythia}. We also benchmark a standard \textbf{SFT} model and a variant, \textbf{Preferred-FT}, which is fine-tuned in a supervised manner on selected completions $y_w$—derived from SFT in the sentiment and summarization experiments, or a generic language model in the single-turn dialogue experiments. Another semi-supervised strategy evaluated is \textbf{Unlikelihood}~\citep{welleck2019neural}, which explicitly trains the model to increase the likelihood of $y_w$ while simultaneously reducing that of $y_l$, modulated by an optional `unlikelihood' coefficient $\alpha\in[0,1]$. We further evaluate \textbf{PPO} \citep{schulman2017proximal}, trained with a reward function derived from preference annotations, as well as \textbf{PPO-GT}, an oracle variant utilizing the ground-truth reward function available in controlled sentiment scenarios. In sentiment experiments, we experiment with two PPO-GT variants: a standard implementation \cite{leandro_von_werra_2023_7790115} and an enhanced version incorporating reward normalization and hyperparameter tuning (these improvements are also adopted when PPO is trained with learned rewards). Additionally, we introduce the \textbf{Best of $N$} approach, which involves sampling $N$ responses from the SFT (or Preferred-FT in the dialogue task) and choosing the response with the highest score according to the reward model trained from preference data. While this strategy can yield strong results by isolating reward model quality from PPO optimization challenges, it is computationally burdensome for practical values of $N$ because it necessitates sampling $N$ completions for each test-time query.

\subsection{How well can DPO optimize the RLHF objective?}

RLHF methods commonly use a KL-constrained reward maximization framework, aiming to maximize reward while keeping the learned policy close to a reference policy. Consequently, when evaluating different algorithms, both the attained reward and the degree of KL divergence must be considered; obtaining marginally higher rewards at the cost of significantly increased KL divergence is generally suboptimal. Figure~\ref{fig:frontier-tldr-main} illustrates the trade-off between reward and KL for several algorithms applied to the sentiment task. For each method, we perform multiple independent training runs, each corresponding to a distinct policy conservativeness parameter (target KL $\in\{3,6,9,12\}$ for PPO, $\beta \in \{0.05,0.1,1,5\}$, $\alpha\in\{0.05,0.1,0.5,1\}$ for unlikelihood, as well as varying random seeds for preferred-FT), for a total of 22 runs. Periodically, after every 100 training iterations and up to convergence, we assess the policies on a designated set of evaluation prompts. For each policy, we measure the mean reward (with respect to the actual reward function) and the mean sequence-level KL divergence\footnote{Specifically, the aggregate over time of the per-step KL-divergences.} to the reference policy, denoted as $\text{KL}\left(\pi\mid \mid \pi_\text{ref}\right)$. The results indicate that DPO establishes a superior frontier, obtaining the best trade-off by achieving higher rewards with lower KL. This finding is noteworthy for several reasons. Not only do DPO and PPO target the same objective, but DPO demonstrates marked efficiency gains, as its reward/KL balance consistently surpasses that of PPO. Furthermore, DPO outperforms PPO even in scenarios where PPO has access to true rewards (PPO-GT).

\subsection{Can DPO scale to real preference datasets?}
\label{sec:dpo-real-datasets}
We proceed to assess how DPO performs when fine-tuned on real-world preference data in the contexts of summarization and single-turn dialogue. For summarization tasks, existing automatic evaluation tools such as ROUGE often fail to align well with human judgments~\citep{stiennon2022learning}, and prior studies have shown that leveraging PPO for fine-tuning on human preferences leads to more effective summaries. Our experimental protocol involves generating completions using different methods on the test portion of the TL;DR summarization dataset and calculating the mean win rate compared to reference outputs in the test set. Completions across all models are drawn at sampling temperatures ranging from 0.0 to 1.0, with win rate results presented in Figure~\ref{fig:frontier-tldr-main} (right). All three methods—DPO, PPO, and Preferred-FT—are built upon fine-tuning the same GPT-J SFT model\footnote{\url{https://huggingface.co/CarperAI/openai_summarize_tldr_sft}}. Our findings show that at a temperature setting of 0.0, DPO attains a win rate of about 61\%, surpassing PPO's peak win rate of approximately 57\% at the same temperature. DPO further outperforms the maximum win rate achieved by the best-of-$N$ baseline. It is important to highlight that $\beta$ for DPO was not systematically optimized, which implies that these results may underrepresent DPO's true capability. Additionally, DPO maintains higher robustness across temperature settings in contrast to PPO, whose effectiveness can drop to the level of the original GPT-J model when the temperature increases. Preferred-FT offers no substantial improvement over the SFT baseline. Direct human evaluations between DPO and PPO, detailed in Section~\ref{sec:human-judgments}, reveal that when sampling with temperature 0.25, DPO completions are favored over PPO completions at temperature 0 by 58\%.

For single-turn dialogue evaluation, we focus on the subset of the Anthropic HH dataset \citep{bai2022training} test split containing interactions limited to a single human-assistant exchange. To assess model performance, we employ GPT-4 to calculate the win rate, referencing the test-preferred completions. In the absence of an established SFT model for this context, we initialize from the pre-trained Pythia-2.8B, train a reference model using Preferred-FT on the selected completions to ensure their compatibility with the model's output distribution, and subsequently apply DPO training. Additionally, we benchmark our models against two baselines: (1) the top completion among 128 generated by Preferred-FT (noting from our investigation that the Best of $N$ metric saturates at $N=128$ for this task; refer to Appendix Figure~\ref{fig:best-of-n}), and (2) a 2-shot prompt variant of the original Pythia-2.8B model. Across these comparisons, DPO consistently matches or outperforms the strongest temperature settings for each method. We further test an RLHF model trained using PPO on the Anthropic HH dataset\footnote{\url{https://huggingface.co/reciprocate/ppo_hh_pythia-6B}}, acquired from a reputable repository\footnote{\url{https://github.com/CarperAI/trlx/tree/main/examples/hh}}, but are unable to identify prompt or sampling temperature configurations yielding improvements over the unmodified Pythia-2.8B. Drawing on findings from TL;DR and given that Best of 128 and PPO both optimize the identical reward signal, we use the Best of 128 as an approximate stand-in for PPO performance. In summary, DPO stands out as the sole efficient approach that surpasses the preferred completions in the Anthropic HH dataset, while also delivering performance comparable to—or exceeding—that of the resource-intensive Best of 128 approach. Notably, Figure~\ref{fig:dialogue-main} illustrates DPO's rapid convergence toward optimal performance.

\subsection{Generalization to a new input distribution}

To further assess how PPO and DPO respond to distributional changes, we apply both policies—trained in the Reddit TL;DR summarization task—to a distinct data distribution: the test split of the CNN/DailyMail news article dataset \citep{nallapati-etal-2016-abstractive}. Sampling temperatures are matched to those found optimal on TL;DR (0 and 0.25). Table~\ref{tab:ood} summarizes the findings. For evaluation, we determine the win rate of GPT-4 over the ground-truth article summaries, employing the same GPT-4 (C) prompt as used in the TL;DR experiments, but adapting the phrase “forum post” to “news article.” On this new domain, DPO continues to substantially exceed PPO in performance. These outcomes suggest that DPO can achieve generalization on par with PPO, despite not leveraging the additional set of unlabeled Reddit TL;DR prompts utilized in PPO training.

\subsection{Validating GPT-4 judgments with human judgments}
\label{sec:human-judgments}
To assess the robustness of GPT-4’s evaluations, we perform a human subject study on the TL;DR summarization task, comparing two different GPT-4 prompt styles. The \textbf{GPT-4 (S)} (simple) prompt requests a selection of the summary that better captures the main content of the post, while the \textbf{GPT-4 (C)} (concise) prompt additionally inquires about conciseness, motivated by our observation that the \textbf{GPT-4 (S)} prompt often favors longer, more repetitive summaries than human raters prefer. Complete prompt details are included in Appendix~\ref{app:prompts}. We conduct three comparative evaluations—using the top-performing model (DPO, temp. 0.25), the lowest-performing model (PPO, temp. 1.0), and an intermediate method (SFT, temp. 0.25)—to encompass a range of summary qualities. Each is compared to PPO with greedy sampling (at its optimal temperature). Results reveal that across both prompt types, GPT-4’s agreement with human raters is comparable to inter-human agreement, indicating GPT-4’s suitability as a proxy for human judgment. Due to limitations on human annotation resources, multiple human ratings are collected solely for the DPO and PPO-1 pairings. In general, the \textbf{GPT-4 (C)} prompt yields win rates more closely aligned with human preferences, prompting us to employ this prompt in the core analyses presented in Section~\ref{sec:dpo-real-datasets}. Additional methodological details, including the rating interface and the roster of human evaluators, are provided in Appendix~\ref{app:human-study}.

\section{Discussion}
Preference-based learning presents a robust and scalable strategy for aligning and improving the capabilities of language models. In this work, we have presented DPO, a straightforward method for training language models from preference data, entirely sidestepping reinforcement learning. Instead of adapting the preference learning objective to a conventional RL formulation to leverage traditional RL tools, DPO establishes a direct correspondence between language model policies and reward representations. This framework enables training models in accordance with human preferences through a cross-entropy loss, omitting both reinforcement learning components and any associated compromises in generality. Remarkably, with minimal hyperparameter adjustment, DPO achieves performance on par with, or surpassing, prevalent RLHF methods such as those utilizing PPO, thereby substantially lowering the practical barriers to scaling up language model training from preference signals.

\textbf{Limitations \& Future Work.} Our findings surface several key areas for future exploration. One question concerns the generalization properties of DPO-trained policies, especially when compared to those developed through explicit reward modeling. Preliminary experiments indicate that DPO-trained policies generalize in ways comparable to models trained with PPO, but a systematic investigation remains necessary. For instance, can self-labeling using the DPO model facilitate efficient use of unlabeled inputs? Another important avenue is to examine how over-optimization of rewards arises in the direct preference optimization paradigm, and whether the modest reduction in performance observed in Figure~\ref{fig:dialogue-main}-right exemplifies this phenomenon. Although our evaluation encompasses models containing up to 6B parameters, scaling DPO to models that are substantially larger represents a promising path for subsequent research. Regarding assessment methodologies, our analysis reveals that GPT-4's win rate metrics are sensitive to the choice of prompt; further research may clarify optimal strategies for eliciting accurate evaluations from automated assessors. Finally, DPO's applications extend beyond preference-based language model training, with potential relevance for developing generative models in additional domains.